\documentclass[a4paper, 12pt]{article}
\usepackage[english]{babel}
\usepackage[utf8]{inputenc}
\usepackage{amstext,amsmath}
\usepackage{parskip}
\usepackage{setspace}
\usepackage{lipsum}  
\usepackage{graphicx}
\usepackage{longtable}
\usepackage{booktabs}
\usepackage{hyperref}
\usepackage{tcolorbox}
\usepackage{csquotes}

\usepackage[
backend=biber,
style=authoryear-comp,
doi=false,isbn=false,url=false,eprint=false,maxcitenames=2,uniquename=false
]{biblatex}

\addbibresource{bibliography.bib} 

\graphicspath{ {images/} }
\onehalfspacing % setspace package allows for 1.5 spacing
% etc. Here you can put all the packages that you will use/nee



\author{Author Name}
\date{\today} 

\begin{document}
These are my first words in LaTeX.

% Table of contents (uncomment to use)
% \tableofcontents

\section{Sectioning}
\subsection{Subsection}
\subsubsection{Subsubsection}

% Unnumbered sections
\section*{Unnumbered Section}
\subsection*{Unnumbered Subsection}
\subsubsection*{Unnumbered Subsubsection}

\section{Inputing Files}
\section{Introduction}
\label{sec:conclusion}

Cite like this: \cite{wooldridgeIntroductoryEconometricsModern2016} or \textcite{wooldridgeIntroductoryEconometricsModern2016}.

\section{Formatting}
% --- Formatting: Font Size and Text Style ---
\textbf{Bold text}, \textit{italic text}, \underline{underlined text}

% Font sizes: \tiny, \scriptsize, \footnotesize, \small, \normalsize, \large, \Large, \LARGE, \huge, \Huge
\begin{itemize}
    \item \tiny tiny
    \item \small small
    \item \normalsize normal
    \item \large large
    \item \Large Large
    \item \LARGE LARGE
    \item \huge huge
    \item \Huge Huge
\end{itemize}

% --- Line Breaks and Spacing ---
First line.\\Second line.\vspace{2.5cm} More vertical space.

% --- Lists ---
\begin{itemize}
    \item First item
    \item Second item
\end{itemize}

\begin{enumerate}
    \item First numbered item
    \item Second numbered item
\end{enumerate}

\section{Mathematical Typesetting}
% --- Mathematical Environments ---
Inline: $x^2 + y^2 = z^2$

Display (unnumbered):
\[
x = \frac{-b \pm \sqrt{b^2-4ac}}{2a}
\]

Display (numbered):
\begin{equation}
    \label{eq:example}
y = mx + c
\end{equation}

Align (multiple lines):
\begin{align}
a &= b + c \\
d &= e - f
\end{align}

% --- Tables ---
\section{Tables}
\label{sec:tables}
Table \ref{tab:example} shows an example of a table.

\begin{table}[h]
    \centering
    \begin{tabular}{lcc}
        \hline
        Variable & Mean & Std. Dev. \\
        \hline
        $x$ & 1.23 & 0.45 \\
        $y$ & 2.34 & 0.67 \\
        \hline
    \end{tabular}
    \caption{Example Table}
    \label{tab:example}
\end{table}


\section{Figures}
\label{sec:figures}
% --- Figures ---
% (Requires an image in the images/ directory)
Figure \ref{fig:example} shows an example of a figure.
\begin{figure}[h]
    \centering
    \includegraphics[width=0.5\textwidth]{logo_uni-2.pdf}
    \caption{Example Figure}
    \label{fig:example}
\end{figure}

% --- Code Listings ---
% (Requires \usepackage{listings})
% \begin{lstlisting}[language=Python, caption=Python example]
% def hello_world():
%     print("Hello, world!")
% \end{lstlisting}

\section{Cross-Referencing}
% --- Cross-Referencing ---
See Table~\ref{tab:example} in section~\ref{sec:tables} for an example table.
See Figure~\ref{fig:example} in section~\ref{sec:figures} for an example figure.
See equation~(\ref{eq:example}) for an example equation.

\section{Citations}
% --- Citations ---
% (Requires biblatex and a .bib file)
Lets cite our favorite paper \cite{ramey2016macroeconomic}. \\
\cite{ramey2016macroeconomic} \\
\textcite{ramey2016macroeconomic} \\
\parencite{ramey2016macroeconomic}  \\
\parencite[p. 23]{ramey2016macroeconomic}  \\
\parencite[e.g.][]{ramey2016macroeconomic}  \\
\citeauthor{ramey2016macroeconomic} \\
\citeyear{ramey2016macroeconomic} \\

% --- Bibliography ---
\appendix
\printbibliography


\end{document}
